\section{The endpoint direction}\label{sec:endpoint}

Theorem~\ref{thm:abelian} sends us transverse to $\omega$, and the proposal's own
phrase --- the variation of the endpoints --- names the direction. The compression
$P_L$ is not in the commutative algebra: the projections commute with one
another, but not with the transfers. This section works out what the
$L$-direction is, and reaches three conclusions, of which the third is the one to
carry forward.

\subsection{The variation is rank one, at the leading endpoint}

Let $V_\omega$ denote the uncompressed causal transfer, so that
$V_{\omega,L}f=P_LV_\omega f$ for $f$ supported in $I_L$. Fix
$f\in C_c^\infty(I_{L_0})$ and let $L>L_0$, so that $f$ is held fixed while the
interval grows.

\begin{proposition}[Endpoint variation]\label{prop:endpoint}
For $L>L_0$ and $f\in C_c^\infty(I_{L_0})$,
\begin{equation}\label{eq:endpoint}
 \frac{\partial}{\partial L}\big(V_{\omega,L}f\big)
 =\tfrac12\,\big(V_\omega f\big)(L/2)\;\delta_{L/2},
\end{equation}
a rank-one variation supported at the single point $x=L/2$, whose coefficient is
the boundary value of the evolved field. Equivalently
$\partial_LV_{\omega,L}=B_LV_{\omega,L}$ with
$B_L=\frac12\,\delta_{L/2}\otimes\mathrm{ev}_{L/2}$.
\end{proposition}

\begin{proof}
$V_{\omega,L}f=\mathbf1_{I_L}\cdot(V_\omega f)$ and
$\partial_L\mathbf1_{(-L/2,L/2)}=\frac12(\delta_{L/2}+\delta_{-L/2})$, so the
variation has the two endpoint terms. Causality kills the second: $f$ is
supported in $I_{L_0}$ and the kernel of $V_\omega$ is supported in
$[0,\infty)$, so $\supp V_\omega f\subseteq[-L_0/2,\infty)$ and
$(V_\omega f)(-L/2)=0$ for $L>L_0$. The surviving coefficient is the boundary
value of $V_{\omega,L}f$ from inside, since $V_\omega f$ and $V_{\omega,L}f$
agree on $I_L$.
\end{proof}

Equation \eqref{eq:endpoint} is exactly the endpoint equation of a line: the
variation is local, it sits at the moving end, and its coefficient is the
transported field there. The tail is fixed and the head moves, which is what
causality means here. This is the closest structural match the proposal has, and
it is not an analogy: the two statements have the same form.

\subsection{The two-parameter connection is flat, so there is no curvature to find}

It is tempting to expect a curvature in the $(\omega,L)$ plane and to look for it
in the mixed derivative. That expectation is wrong, and it is worth saying why,
because the correction is what sends the search somewhere else.

\begin{proposition}[Flatness]\label{prop:flat}
On the range of the two-parameter family, wherever $V_{\omega,L}f$ is twice
continuously differentiable in $(\omega,L)$,
\begin{equation}\label{eq:flat}
 \partial_\omega B_L+\partial_LG_{\omega,L}+[G_{\omega,L},B_L]=0 .
\end{equation}
\end{proposition}

\begin{proof}
Mixed partial derivatives of $V_{\omega,L}f$ commute. Substituting
$\partial_\omega V_{\omega,L}=-G_{\omega,L}V_{\omega,L}$ from \eqref{eq:flow} and
$\partial_LV_{\omega,L}=B_LV_{\omega,L}$ from Proposition~\ref{prop:endpoint}
into $\partial_\omega\partial_L=\partial_L\partial_\omega$ and collecting terms
gives \eqref{eq:flat} applied to $V_{\omega,L}f$.
\end{proof}

So the pair $(-G_{\omega,L}\dd\omega,\,B_L\dd L)$ is a \emph{flat} connection
wherever it is a connection at all: the zero-curvature condition is not an extra
requirement to be imposed, it is forced by the existence of the two-parameter
family. There is no local curvature to compute and no surface term to extract,
and the non-Abelian Stokes theorem still has nothing to integrate. On a simply
connected parameter domain, the holonomy of a flat connection depends on the
endpoints alone.

\subsection{The connection does not close, and it misses by exactly the borderline}

Proposition~\ref{prop:endpoint} holds on the core $C_c^\infty(I_L)$, where
$V_\omega f$ is smooth because $\widehat f$ decays rapidly. It does not extend to
the form domain.

\begin{proposition}[No form-bounded endpoint evaluation]\label{prop:trace}
$B_L$ is not bounded relative to $Q_{0,L}$, at any $0\leq\omega\leq\frac12$.
\end{proposition}

\begin{proof}[Proof sketch]
$B_L$ is built from $\mathrm{ev}_{L/2}$, and a point evaluation on $\R$ is
bounded on $H^r$ only for $r>\frac12$. The form domain is
$\mathcal D_{\log,L}$ of Definition~\ref{def:generator}, whose multiplier
$b(\tau^2)=\log|\tau|+O(1)$ is only logarithmically stronger than $L^2$; it is
contained in no $H^r$ with $r>0$ and carries no trace operator, which is also why
the natural closed domain of the archimedean energy has no trace condition
\cite{InverseBulk}. The transfer does not repair this: the bound
$|K_\omega(\eta+i\tau)|\leq C_\eta(2+|\tau|)^{-\omega}$ of \cite{Background}
gives $V_\omega$ a smoothing gain into $H^r$ for $0<r<\min(\omega,\frac12)$ only,
which is strictly below the trace threshold for every admissible shift, the
endpoint $\omega=\frac12$ included.
\end{proof}

\begin{remark}[A third marginality at the same exponent]\label{rem:marginal}
The failure in Proposition~\ref{prop:trace} is not by a wide margin; it is
exactly borderline. Three separate quantities in this program sit at the same
exponent and on the wrong side of it: the archimedean kernel
$\gammak(r)\sim\frac1{2r}$ makes $K$ the Gagliardo seminorm at $s=0$, the
borderline of the scale, so the form domain is only logarithmically better than
$L^2$; the transfer gains at most $(2+|\tau|)^{-1/2}$; and a trace needs
$r>\frac12$. We record the coincidence and do not explain it. Whether it is one
fact seen three times is Problem~\ref{prob:trace}.
\end{remark}

\subsection{The primes are the atoms of the connection}

The proof of Theorem~\ref{thm:kernel} decomposes the kernel of $G_{\omega,L}$
into a continuous part and an atomic part:
\begin{equation}\label{eq:atomic}
\begin{aligned}
 g_\omega(t)&=g_\omega^{\rm cont}(t)+g_\omega^{\rm at}(t),\qquad t>0,\\
 g_\omega^{\rm cont}(t)&=4\cosh(\omega t)\cosh(t/2)
   -2\cosh(\omega t)\gammak(t)+\mathrm{const},\\
 g_\omega^{\rm at}(t)&=-2\sum_{n\geq2}\Lambda(n)n^{-1/2}
   \cosh(\omega\log n)\,\delta_{\log n},
\end{aligned}
\end{equation}
the atoms sitting at the prime periods $m\log p$ with masses
$2(\log p)p^{-m/2}\cosh(\omega m\log p)$. This is the L\'evy decomposition that
the companion investigation's ninth selection rule requires of any source
\cite{Selection}, and it is here a property of the connection rather than of a
prepared state.

Read along the line, an atomic connection is a connection with \emph{punctures}:
the holonomy of the segment is an ordered product of one local factor per atom
crossed, times the contribution of the continuous part. That is precisely the
periodic-orbit structure the program's own analysis says the primes demand and
which a short-distance mechanism cannot supply \cite{Bottleneck}. It is also,
given Propositions~\ref{prop:flat} and \ref{prop:trace}, the only place left for
the proposal's mechanism to live: a flat connection has no local curvature, but a
flat connection on a punctured domain has monodromy, and monodromy is exactly the
non-local, potentially non-abelian datum the proposal needs.

The obstruction is stated by Theorem~\ref{thm:abelian}. The local factors at
different atoms commute, because every object here is a convolution and the
atomic masses are \emph{scalars}. The monodromy representation is therefore
one-dimensional and carries no more information than the product of its masses.
To make it carry more, the residues at the atoms would have to be matrix valued
--- and the arithmetic object whose prime sum carries matrix-valued residues is
$\operatorname{tr}\rho(\mathrm{Frob}_p)$ for a representation $\rho$. The three
strands of Remark~\ref{rem:cost} converge on this point: the endpoint direction
is where the geometry is, the atoms are where the primes are, and an internal
index at the atoms is the only way to make the two non-abelian at once. We state
this as the structural conclusion of the paper and not as a programme, since it
changes the target from $\zeta$ to a family.

\subsection{Two open problems}

\begin{problem}[Close the endpoint flow, or show it cannot be closed]
\label{prob:trace}
Proposition~\ref{prop:trace} obstructs $B_L$ on the form domain while
Proposition~\ref{prop:endpoint} defines it on a dense core. Determine whether the
endpoint flow admits a closed realization on any domain on which
$Q_{\omega,L}$ is also closed --- for instance by pairing the endpoint evaluation
against the atomic part of \eqref{eq:atomic}, which is where the flow's own data
sits --- or prove that the borderline failure of Remark~\ref{rem:marginal} is
structural.
\end{problem}

\begin{problem}[Monodromy of the punctured flow]\label{prob:curvature}
Compute the local factor of \eqref{eq:atomic} at a single atom, as an operator on
the core, and determine what a matrix-valued residue at that atom would have to
satisfy for the resulting holonomy to remain a contraction in the sense of
Proposition~\ref{prop:route}. This is the smallest calculation that would tell
whether the internal index of Remark~\ref{rem:cost}(3) can be introduced without
destroying the flow identities.
\end{problem}

\begin{remark}[What the $L$-direction is not]\label{rem:not-a-family}
Two cautions, both of which this program has paid for once. First, $Q_L$ is not a
family of functionals but one functional restricted by support, so a realization
at fixed $L$ carries no information and an $L$-family of unrelated realizations
proves nothing \cite{InverseBulk}. The $L$-dependence that is real is in the
compression, which is what Proposition~\ref{prop:endpoint} differentiates.
Second, the exclusions this program has accumulated are statements about channel
decompositions; a flow is not a decomposition, so they neither block this
direction nor may be quoted in its support \cite{ExclusionMap}.
\end{remark}
